\scanpage{71}
triangle de triades déliées (parallèles) la contenant, sont formé de triades contenues dans $S\setminus X$.

7) Considérons l'application canonique surj.

\[
S\setminus X=X_1\amalg X_2\xrightarrow{\pi}I_0
\]

dont les fibres sont les réunions $\tau_1\cup\tau_2=\tau_1\cup\overline{\tau}_1$, où $\tau_1\inserted{\unsure{\in|e}}$ est une triade de $X_1$ de la famille $I'_1$, et $\tau_2$ \deleted{\unsure{la}} \inserted{une} triade associée de $X_2$ \inserted{de la famille $I'_2$} (donc $\tau_2=\overline{\tau}_1$). Alors

\unsure{a)} Pour tt $i\in I_0$, la réunion $X\cup\pi^{-1}(\{i\})$ \deleted{est} un $\unsure{\mathscr T|\Gamma|\mathcal T}_1$ plongé dans $S$, \inserted{et contenant $X$} et l'application $i\longmapsto X\cup\pi^{-1}(\{i\})$ est une bijection de $I_0$ sur l'ens des \deleted{\unsure{3} $\mathscr T_1$} tels $\mathscr T_1$ \inserted{$\unsure{\mathscr X|\mathfrak X}$}. Donc les $\pi^{-1}(i)$ sont les $\mathscr X\setminus X$, où $\mathscr X$ est un $\mathscr T_1$ de $S$ contenant $X$.

8) Ces ens sont aussi les orbites dans $S\setminus X$ sous le sous-gpe de $\mathrm{Aut}(X)$ engendré par les symétries \deleted{\unsure{par aux} triangles} \inserted{$\unsure{\sigma_p|\sigma_t}$, $\unsure{\in|\underline t}$} de $X$.

[\unsure{En|Bien} fait, le gpe des autom ``int'' de $X$ opère sur $S$ \ldots]

% archive sheet number: 70

